\section{Derivations}\label{app:derivations}
This appendix derives the small results that Chapter~\ref{chp:method_cell} uses without proof.

\subsection{Depth Error of a Stereo Camera}
A stereo camera with focal length $f$ (pixels) and baseline $b$ measures depth from the disparity
$d_{\mathrm{px}}$ as $z=fb/d_{\mathrm{px}}$. A disparity error $\delta$ changes the depth, to first order, by
\begin{equation}
\Delta z\approx-\frac{fb}{d_{\mathrm{px}}^2}\,\delta=-\frac{z^2}{fb}\,\delta ,
\end{equation}
whose magnitude is Equation~\ref{eq:depth_error}; an overestimated disparity shortens the depth. A
constant disparity offset therefore produces an error that grows with the square of the range, and a
check at one range cannot separate it from a constant depth offset: two ranges determine $a$ and $k$ in
$\Delta z=a+kz^2$, and a third tests the quadratic form against a linear one.

\subsection{Tool-Centre Error from Touches on a Plane}
Let the plane be $\vect{n}^{\top}\vect{x}+d=0$, measured with the TCP in use. A touch records the
tool pose $(\mathbf{R}_i,\vect{p}_i)$ at the TCP; if the real tip is off by $\vect{e}$ in the tool
frame, it lies at $\vect{p}_i+\mathbf{R}_i\vect{e}$, on the plane:
\begin{equation}
\vect{n}^{\top}(\vect{p}_i+\mathbf{R}_i\vect{e})+d=0
\quad\Longrightarrow\quad
(\mathbf{R}_i^{\top}\vect{n})^{\top}\vect{e}=-h_i,\qquad h_i=\vect{n}^{\top}\vect{p}_i+d .
\end{equation}
Since the plane itself was measured with the erroneous TCP, its offset $d$ is biased by the vertical
component of the error, which a free offset $\delta$ absorbs, giving Equation~\ref{eq:tcp_check}.
With the pen vertical (the tool's $z$ axis pointing down to the tip, $\vect{n}$ up),
$\mathbf{R}_i^{\top}\vect{n}\approx-\vect{e}_z$ for every roll, so only $e_z$
(confounded with $\delta$) is seen; tilting the pen by $\vartheta$ towards azimuth $\phi$ makes
$\mathbf{R}_i^{\top}\vect{n}$ gain a lateral component of size $\sin\vartheta$ in the direction
$\phi$. Tilts towards four azimuths therefore make $e_x$ and $e_y$ observable with a sensitivity of
$\sin\vartheta$, and the along-pen component separates from $\delta$ only through
$1-\cos\vartheta$, which is small at $30^\circ$ (0.13) and better at $45^\circ$ (0.29).

\subsection{Phase Margin of the Camera Aim}
For one axis, the error angle $e$ obeys $\dot\theta_c=K\,e(t-\tau)$, with $e=\theta_t-\theta_c$,
a lag $\tau$ from the smoother and the arm, and the latency of the image compensated
(\S\ref{sec:mc_aim}). The open-loop transfer function is $L(s)=K\,e^{-s\tau}/s$, so
\begin{equation}
|L(j\omega)|=\frac{K}{\omega},\qquad \angle L(j\omega)=-\frac{\pi}{2}-\omega\tau .
\end{equation}
The gain crosses over at $\omega_c=K$, and the phase margin is
$\pi+\angle L(j\omega_c)=\pi/2-K\tau$, Equation~\ref{eq:aim_pm}. A first-order lag of time constant
$\tau_s$ contributes $\arctan(\omega\tau_s)\approx\omega\tau_s$ of phase at low frequency, so the two
60~ms lags of the smoother count as about 0.12~s of the lag.

\subsection{Latency Compensation by Rotation About the Optical Centre}
If the camera rotates about its optical centre $\vect{c}$, fixed in the base frame, and the fingertip
$\vect{x}$ is still, the unit direction $\vect{r}_B=(\vect{x}-\vect{c})/\norm{\vect{x}-\vect{c}}$ is
constant. A ray measured in the camera frame at time $t_k$, $\hat{\vect{d}}_k$, satisfies
$\vect{r}_B=\Rf{B}{C}(t_k)\,\hat{\vect{d}}_k$; at any later time the same direction in the current
camera frame is $\Rf{B}{C}(t)^{\top}\vect{r}_B$. The error computed from an old image is therefore the
exact current error, independent of $\norm{\vect{x}-\vect{c}}$. If the centre moves by
$\Delta\vect{c}$, the true direction changes by the angle between $\vect{x}-\vect{c}$ and
$\vect{x}-\vect{c}-\Delta\vect{c}$, at most $\arcsin(\norm{\Delta\vect{c}}/\norm{\vect{x}-\vect{c}})$:
with the optical-centre hold keeping $\norm{\Delta\vect{c}}$ to millimetres at half a metre, the
residual is a fraction of a degree.

\subsection{Information Left for the Extrinsic Translation}
At its solution, the joint estimate of Equation~\ref{eq:extrinsic_d} is a linearised, iteratively
reweighted least-squares problem, with its rows weighted by the Welsch and view weights, in the
unknowns $\vect{x}$, the part corrections, and $\vect{d}$. Its normal matrix is
\begin{equation}
\mathbf{N}=\begin{bmatrix}\mathbf{N}_{xx}&\mathbf{N}_{xd}\\ \mathbf{N}_{dx}&\mathbf{N}_{dd}\end{bmatrix},
\qquad
\operatorname{Cov}(\vect{d})=\sigma^2\big(\mathbf{N}^{-1}\big)_{dd}=\sigma^2\big(\mathbf{N}_{dd}-\mathbf{N}_{dx}\mathbf{N}_{xx}^{-1}\mathbf{N}_{xd}\big)^{-1},
\end{equation}
by block inversion. The Schur complement $\mathbf{S}$, divided by $\sigma^2$, is therefore the
information about $\vect{d}$ after the poses have absorbed what they can: a shift that a change of the part poses explains as well
as $\vect{d}$ does contributes nothing to $\mathbf{S}$. A pseudo-inverse replaces $\mathbf{N}_{xx}^{-1}$
when a part has an unobservable direction (a tube's spin).

\subsection{The Overshoot Clearance Bound}
During the descent the tool translates along its own axis $\vect{a}$ with its orientation held. For
any point $\vect{y}$ of the tool and any obstacle $\mathcal{O}$, the distance function
$\operatorname{dist}(\cdot,\mathcal{O})$ is 1-Lipschitz, so
\begin{equation}
\operatorname{dist}(\vect{y}+s\,\vect{a},\mathcal{O})\ \ge\ \operatorname{dist}(\vect{y},\mathcal{O})-s .
\end{equation}
Taking the minimum over the tool's points, a pose $s$ past the reference pose keeps at least the
reference clearance minus $s$. For the tool's capsules, whose distances to boxes and to the table are
true distances, the bound can therefore never be violated. What the descent check can refuse lies
outside the argument: pairs of boxes, which are only tested for intersection (such as the camera body
against a part box), and the arm's capsules, which do not translate rigidly.

\subsection{Clearance of the Pen Holder in a Fillet}
Let a fillet have opening $\theta$, the pen lie on its bisector with the tip at the root, and a
capsule of radius $r_h$ have its axis along the pen from $\ell$ to $\ell'>\ell$ above the tip. A
point at height $h$ on the bisector is at distance $h\sin(\theta/2)$ from each plate, increasing
with $h$, so the capsule's nearest point to the plates is on the sphere at $h=\ell$, at a distance
$\ell\sin(\theta/2)-r_h$: Equation~\ref{eq:holder_clearance}. For a work-angle tilt $\vartheta$ away
from one plate (with $\theta/2+\vartheta\le90^\circ$), the clearance to that plate grows to
$\ell\sin(\theta/2+\vartheta)-r_h$ and the clearance to the other shrinks to
$\ell\sin(\theta/2-\vartheta)-r_h$. With the pen on the bisector both are equal and the smaller is
largest at $\vartheta=0$; a tilt helps only when one side binds, as with a leaning plate or a tool that
is not symmetric about the pen.

\subsection{Six Boxes Round a Tube}
Each of the six boxes spans the tube's diameter $2r$ in one direction and $2w$, $w=r\sin15^\circ$, in
the other, rotated by multiples of $30^\circ$. A point of the disc at radius $\rho\le r$ and angle
$\phi$ is within $15^\circ$ of some box axis, so its distance from that axis is
$\rho\,|\sin(\phi-\phi_j)|\le r\sin15^\circ=w$, and its distance along the axis,
$\rho\,|\cos(\phi-\phi_j)|$, is at most $r$; the box covers it: the union covers the disc. The
union's farthest points are the box corners, at $\sqrt{r^2+w^2}$ from the axis, so the union stands at
most
\begin{equation}
r\Big(\sqrt{1+\sin^2 15^\circ}-1\Big)\approx0.033\,r
\end{equation}
proud of the wall, against $r(\sqrt2-1)\approx0.41\,r$ for a single box round the tube.
